\chapter{Complexity — Universal Interpreter}
%%% generated by the blueprint pipeline from TCSlib.Complexity.TuringMachine.UniversalInterpreter
%%%%%%%%%%%%%%%%%%%%%%%%%%%%%%%%%%%%%%%%%%%%%%%%%%%%%%%%%%%%%%%%%%%%%%%%%%%%%%%
\section{Overview}

This module builds the table interpreter at the heart of the universal Turing
machine (Arora--Barak, \S 1.4.1, Theorem~1.9): a fixed four-tape machine over
$\{0,1\}$ whose work tapes hold a serialized transition table, a unary state
counter, the simulated work tape, and a virtual-left-boundary marker.  It
provides the interpreter's finite control, the complete universal-machine
candidate, exact-cost gadget lemmas for every table and state-tape scanning
phase, exact interpreter initialization from a canonical serialization, and
the assembly of both evaluator clauses from a configuration-level block
simulation.  Throughout, bit strings are written over $\{0,1\}$ ($1$ for true,
$0$ for false); the \emph{pair encoding} $\langle \alpha, x\rangle$ of two bit
strings lists each bit of $\alpha$ twice, then the unequal delimiter pair
$01$, then $x$ verbatim; and the \emph{buffer tape} of a string $w$ is the
two-way infinite tape holding $w$ in cells $0,\dots,|w|-1$ and blanks
elsewhere.

%%%%%%%%%%%%%%%%%%%%%%%%%%%%%%%%%%%%%%%%%%%%%%%%%%%%%%%%%%%%%%%%%%%%%%%%%%%%%%%
\section{Declarations}

\begin{definition}[Finite control of the table interpreter]
\lean{Turing.UniversalControl}
\label{Turing.UniversalControl}
\leanok
The finite control-state type of the universal machine's table interpreter: a
fixed collection of phases (table rewind, count parsing, initial-state copying
and skipping, state rewind, main dispatch, record skipping and reading,
next-state copying, record application), each carrying only bounded registers
--- Boolean flags, record indices below $9$, a field index below $8$, and a
buffer of eight action bits --- together with a sink state for malformed
reads.  In particular the control is finite and independent of the number of
states of the simulated machine; the unbounded simulated state is represented
only on a tape.
\end{definition}

\begin{definition}[Quadruple of tape entries]
\lean{Turing.universalFour}
\label{Turing.universalFour}
\leanok
Given four values $a, b, c, d$ of a common type, the function on the
four-element index type sending index $0$ to $a$, index $1$ to $b$, index $2$
to $c$, and index $3$ to $d$.  It packages one entry per work tape of a
four-tape machine into a single function, without a variable-size register.
\end{definition}

\begin{definition}[Administrative interpreter action]
\lean{Turing.universalAdmin}
\label{Turing.universalAdmin}
\leanok
\uses{Turing.Action, Turing.UniversalControl, Turing.universalFour}
For a target control state $q$, a table-cursor movement
$d_t \in \{-1, 0, +1\}$, and a state-tape instruction consisting of an
optional write and a cursor movement $d_s$, the four-tape machine action that
keeps the input head fixed, writes nothing and stays put on the
simulated-work and marker tapes, moves the table cursor by $d_t$, performs
the optional write and the movement $d_s$ on the state tape, emits no output,
and passes to control state $q$.  By default $d_t = 0$ and the state tape is
untouched.
\end{definition}

\begin{definition}[Serialization index of a read symbol]
\lean{Turing.universalReadIndex}
\label{Turing.universalReadIndex}
\leanok
The position of a tape symbol in the fixed serialization order of the three
possible reads: the blank has index $0$, the bit $0$ has index $1$, and the
bit $1$ has index $2$.
\end{definition}

\begin{definition}[Record index within a state block]
\lean{Turing.universalRecordIndex}
\label{Turing.universalRecordIndex}
\leanok
\uses{Turing.universalReadIndex}
For a pair of read symbols --- an input symbol and a work symbol, each blank,
$0$, or $1$ --- the index in $\{0, \dots, 8\}$ of the corresponding transition
record within the nine-record block of one coded state: three times the
serialization index of the input read (blank $\mapsto 0$, bit $0 \mapsto 1$,
bit $1 \mapsto 2$) plus the serialization index of the work read.
\end{definition}

\begin{definition}[Decoding the two-bit movement field]
\lean{Turing.universalSign}
\label{Turing.universalSign}
\leanok
Decodes a valid two-bit head-movement field into a movement in
$\{-1, 0, +1\}$: the bits $11$ decode to $-1$ (left), the bits $10$ decode to
$+1$ (right), and a first bit $0$ decodes to $0$ (stay).
\end{definition}

\begin{definition}[Decoding the two-bit optional-write field]
\lean{Turing.universalWrite}
\label{Turing.universalWrite}
\leanok
Decodes a valid two-bit optional-write field: the bits $1b$ mean ``write the
bit $b$'', the bits $01$ mean ``write the blank'' (erase the cell), and the
bits $00$ mean ``leave the cell unchanged''.
\end{definition}

\begin{definition}[The four-tape table interpreter]
\lean{Turing.universalInterpreter}
\label{Turing.universalInterpreter}
\leanok
\uses{Turing.FinTM.virtualMove, Turing.MultiTapeTM, Turing.UniversalControl,
  Turing.universalAdmin, Turing.universalFour, Turing.universalRecordIndex,
  Turing.universalSign, Turing.universalWrite}
The fixed four-tape machine over $\{0,1\}$ that interprets a serialized
transition table.  Its four work tapes hold the captured program table, the
current simulated state in unary behind a permanent origin marker, the
simulated machine's work tape, and a marker tape recording whether the
virtual input head sits at the simulated left boundary.  To perform one
simulated transition it rewinds the table, skips nine records per unary state
symbol while consuming the unary state destructively, selects within the
reached block the record indexed by the current simulated reads, buffers its
eight fixed action bits, copies the successor-state unary field back onto the
state tape, and applies the decoded action, adjusting the input-head movement
at the virtual left boundary.  The control is finite and independent of the
number of coded states; malformed administrative reads enter a non-halting
sink state.
\end{definition}

\begin{definition}[Unary state-tape representation]
\lean{Turing.universalStateTape}
\label{Turing.universalStateTape}
\leanok
\uses{Turing.FinTM.bufferTape}
The tape encoding of a simulated state index $q \in \mathbb{N}$: cell $0$
holds the permanent origin marker $0$, cells $1, \dots, q$ hold $1$s (the
unary representation of $q$), and all other cells are blank.
\end{definition}

\begin{definition}[The complete universal machine candidate]
\lean{Turing.universalTM}
\label{Turing.universalTM}
\leanok
\uses{Turing.EffectiveMachineCode, Turing.FinTM, Turing.universalCanonTM,
  Turing.universalCaptureTM, Turing.universalInterpreter}
The universal-machine candidate associated with an effective machine-coding
scheme: the composite finite-tape machine over $\{0,1\}$ that first virtually
runs the scheme's canonizer to extract the code prefix of the pair-encoded
input and produce the canonical serialized table, then captures that table
onto a work tape, and finally hands control to the fixed four-tape table
interpreter.
\end{definition}

\begin{definition}[Interpreter evaluation configuration]
\lean{Turing.universalEvalCfg}
\label{Turing.universalEvalCfg}
\leanok
\uses{Turing.Cfg, Turing.FinTM.bufferTape, Turing.UniversalControl,
  Turing.universalFour}
Given a base four-tape configuration (supplying the physical input position,
the simulated work tape, the boundary-marker tape, and the accumulated
output), a control state $q$, a bit string $T$ with a table cursor
$p \in \mathbb{Z}$, and a state tape
$\sigma \colon \mathbb{Z} \to \{0, 1, \mathrm{blank}\}$ with a state cursor
$s$: the configuration in control state $q$ whose first work tape is the
buffer tape of $T$ with head at $p$, whose second work tape is $\sigma$ with
head at $s$, and which agrees with the base configuration in every other
field.
\end{definition}

\begin{lemma}[Reads of an evaluation configuration]
\lean{Turing.universalEval_reads}
\label{Turing.universalEval_reads}
\leanok
\uses{Turing.Cfg, Turing.Cfg.workTapeSymbols, Turing.FinTM.bufferTape,
  Turing.UniversalControl, Turing.universalEvalCfg, Turing.universalFour}
\difficulty{3}
Consider an interpreter evaluation configuration whose table tape is the
buffer tape of a bit string $T$ with cursor $p$, whose state tape is
$\sigma$ with cursor $s$, and whose remaining data come from a base
configuration.  The four symbols under its work heads are, in order: the
symbol of the buffer tape of $T$ at $p$, the state symbol $\sigma(s)$, and
the two symbols read by the base configuration on its third and fourth work
tapes.
\end{lemma}

\begin{lemma}[Effect of one administrative action]
\lean{Turing.universalAdmin_apply}
\label{Turing.universalAdmin_apply}
\leanok
\uses{Turing.Action.apply, Turing.Cfg, Turing.UniversalControl,
  Turing.moveInputPos_zero, Turing.universalAdmin, Turing.universalEvalCfg,
  Turing.universalStateWindow}
\difficulty{3}
Applying the administrative action with target control state $q'$, table
movement $d_t$, optional state-tape write $w$, and state movement $d_s$ to an
interpreter evaluation configuration with table cursor $p$ and state tape
$\sigma$ with cursor $s$ yields the evaluation configuration over the same
base with control state $q'$, table cursor $p + d_t$, state tape equal to
$\sigma$ updated at $s$ by $w$ when a write is present (and $\sigma$
otherwise), and state cursor $s + d_s$.  The input head, the other two work
tapes, and the accumulated output are unchanged.
\end{lemma}

\begin{lemma}[Administrative step rule]
\lean{Turing.universalEval_step}
\label{Turing.universalEval_step}
\leanok
\uses{Turing.Action.apply, Turing.Cfg, Turing.Cfg.inputSymbol,
  Turing.Cfg.workTapeSymbols, Turing.FinTM.bufferTape,
  Turing.MultiTapeTM.step, Turing.UniversalControl, Turing.universalAdmin,
  Turing.universalAdmin_apply, Turing.universalEvalCfg,
  Turing.universalEval_reads, Turing.universalFour,
  Turing.universalInterpreter}
\difficulty{3}
Consider an evaluation configuration of the table interpreter with control
state $q$, table $T$ with cursor $p$, and state tape $\sigma$ with cursor $s$
over some base configuration.  If the interpreter's transition function at
$q$, applied to the physical input symbol and the represented reads (the
buffer-tape symbol of $T$ at $p$, the state symbol $\sigma(s)$, and the base
configuration's two remaining work reads), is the administrative action with
target $q'$, table movement $d_t$, optional state write $w$, and state
movement $d_s$, then one interpreter step takes the configuration to the
evaluation configuration with control state $q'$, table cursor $p + d_t$,
state tape updated by $w$ at $s$, and state cursor $s + d_s$, leaving all
other fields unchanged.
\end{lemma}

\begin{lemma}[Reading the first unconsumed table cell]
\lean{Turing.universal_table_read}
\label{Turing.universal_table_read}
\leanok
\uses{Turing.FinTM.bufferTape, Turing.FinTM.bufferTape_nat}
\difficulty{2}
For bit strings $l, r$ and a bit $b$, the buffer tape of the concatenation
$l \cdot b \cdot r$ (holding the string in cells $0, 1, \dots$ with blanks
elsewhere) reads exactly $b$ at position $|l|$.
\end{lemma}

\begin{lemma}[Exact-cost table rewind]
\lean{Turing.universal_table_rewind}
\label{Turing.universal_table_rewind}
\leanok
\uses{Turing.Cfg, Turing.FinTM.bufferTape, Turing.FinTM.bufferTape_nat,
  Turing.MultiTapeTM.runFrom, Turing.MultiTapeTM.runFrom_succ_eq_step,
  Turing.MultiTapeTM.runFrom_zero, Turing.UniversalControl,
  Turing.universalEvalCfg, Turing.universalEval_step, Turing.universalFour,
  Turing.universalInterpreter}
\difficulty{4}
Consider the table interpreter in its table-rewind phase (carrying a Boolean
flag and a record index that it merely preserves), in an evaluation
configuration over a table $T$ with table cursor $j - 1$ for some
$j \le |T|$, and with an arbitrary state tape $\sigma$ and state cursor $s$.
Then exactly $j + 1$ interpreter steps bring it to the count-parsing phase
with the same flag and index and with the table cursor at $0$; the state
tape, state cursor, and all other data are unchanged.
\end{lemma}

\begin{lemma}[Exact-cost count-field skip]
\lean{Turing.universal_count_run}
\label{Turing.universal_count_run}
\leanok
\uses{Turing.Cfg, Turing.FinTM.bufferTape, Turing.MultiTapeTM.runFrom,
  Turing.MultiTapeTM.runFrom_add, Turing.MultiTapeTM.runFrom_succ_eq_step',
  Turing.MultiTapeTM.runFrom_zero, Turing.MultiTapeTM.step,
  Turing.UniversalControl, Turing.universalEvalCfg,
  Turing.universalEval_step, Turing.universalFour,
  Turing.universalInterpreter, Turing.universalStateWindow,
  Turing.universal_table_read}
\difficulty{5}
Suppose the table is $l \cdot d(\beta) \cdot 01 \cdot r$, where $d(\beta)$
doubles each bit of a bit string $\beta$ and $01$ is the unequal delimiter
pair.  Starting the interpreter's count parser (carrying a Boolean flag and a
record index) at table cursor $|l|$, exactly $2|\beta| + 2$ steps move the
table cursor past the doubled count field and the delimiter to position
$|l| + 2|\beta| + 2$, entering the initial-state copier when the flag is set
and the initial-state skipper with the carried index otherwise.  No binary
arithmetic on the value of $\beta$ is performed, and the state tape, state
cursor, and all other data are unchanged.
\end{lemma}

\begin{definition}[State window during destructive lookup]
\lean{Turing.universalStateWindow}
\label{Turing.universalStateWindow}
\leanok
The state-tape contents during a destructive table lookup, parametrized by
the number $c$ of unary cells already erased and the number $m$ still
remaining: cell $0$ holds the permanent origin marker $0$, cells
$c + 1, \dots, c + m$ hold $1$s, and all other cells are blank.
\end{definition}

\begin{lemma}[The intact window is the state tape]
\lean{Turing.universalStateWindow_zero}
\label{Turing.universalStateWindow_zero}
\leanok
\uses{Turing.FinTM.bufferTape, Turing.universalStateTape,
  Turing.universalStateWindow}
\difficulty{4}
With no cells erased, the state window with $n$ remaining unary cells ---
marker $0$ at the origin, $1$s in cells $1, \dots, n$, blanks elsewhere ---
coincides with the intact unary state tape encoding the state index $n$.
\end{lemma}

\begin{lemma}[A nonempty window reads one at its cursor]
\lean{Turing.universalStateWindow_read}
\label{Turing.universalStateWindow_read}
\leanok
\uses{Turing.universalStateWindow}
\difficulty{2}
In a state window with $j$ erased cells and $n + 1$ remaining unary cells,
the cell at position $j + 1$ (the current lookup position) holds the bit $1$.
\end{lemma}

\begin{lemma}[An exhausted window reads blank at its cursor]
\lean{Turing.universalStateWindow_end}
\label{Turing.universalStateWindow_end}
\leanok
\uses{Turing.universalStateWindow}
\difficulty{2}
In a state window with $j$ erased cells and no remaining unary cells, the
cell at position $j + 1$ is blank.
\end{lemma}

\begin{lemma}[Erasing one unary symbol]
\lean{Turing.universalStateWindow_erase}
\label{Turing.universalStateWindow_erase}
\leanok
\uses{Turing.universalStateWindow}
\difficulty{4}
Erasing the cell at position $j + 1$ of a state window with $j$ erased and
$n + 1$ remaining unary cells (writing the blank there) yields exactly the
state window with $j + 1$ erased and $n$ remaining cells.
\end{lemma}

\begin{lemma}[A fully consumed window is marker-only]
\lean{Turing.universalStateWindow_empty}
\label{Turing.universalStateWindow_empty}
\leanok
\uses{Turing.universalStateWindow}
\difficulty{3}
A state window with no remaining unary cells is the same tape regardless of
the number of erased cells: the permanent origin marker $0$ at position $0$
and blanks everywhere else.
\end{lemma}

\begin{lemma}[Appending a unary state symbol]
\lean{Turing.universalStateTape_append}
\label{Turing.universalStateTape_append}
\leanok
\uses{Turing.FinTM.bufferTape_append, Turing.universalStateTape}
\difficulty{2}
Writing the bit $1$ into cell $n + 1$ of the intact unary state tape encoding
the index $n$ yields the intact unary state tape encoding the index $n + 1$.
\end{lemma}

\begin{lemma}[The intact state tape ends in blank]
\lean{Turing.universalStateTape_end}
\label{Turing.universalStateTape_end}
\leanok
\uses{Turing.FinTM.bufferTape, Turing.universalStateTape}
\difficulty{1}
The intact unary state tape encoding the index $n$ is blank at position
$n + 1$, immediately after its last unary symbol.
\end{lemma}

\begin{lemma}[Exact-cost marker-directed state rewind]
\lean{Turing.universal_state_rewind}
\label{Turing.universal_state_rewind}
\leanok
\uses{Turing.Cfg, Turing.MultiTapeTM.runFrom,
  Turing.MultiTapeTM.runFrom_succ_eq_step, Turing.MultiTapeTM.runFrom_zero,
  Turing.UniversalControl, Turing.universalAdmin, Turing.universalEvalCfg,
  Turing.universalEval_step, Turing.universalFour,
  Turing.universalInterpreter}
\difficulty{4}
Let $\sigma$ be a state tape whose cell $0$ reads the marker bit $0$ and none
of whose cells at positive natural positions reads $0$ (they may be $1$ or
blank).  Let $q$ be an interpreter control state such that, for every
combination of reads whose state-tape read is $0$, the transition at $q$ is
the administrative action passing to a fixed state $q'$ and moving the state
cursor right, and for every combination whose state-tape read is not $0$ it
is the administrative action staying in $q$ and moving the state cursor left,
in both cases touching nothing else.  Then for every natural number $j$,
starting from the evaluation configuration in state $q$ with state tape
$\sigma$ and state cursor $j$, exactly $j + 1$ interpreter steps reach state
$q'$ with state cursor $1$, leaving $\sigma$, the table cursor, and all other
data unchanged.  The hypotheses are deliberately independent of whether the
traversed cells are erased blanks or retained unary symbols.
\end{lemma}

\begin{lemma}[Exact-cost initial-state skip]
\lean{Turing.universal_initial_skip}
\label{Turing.universal_initial_skip}
\leanok
\uses{Turing.Cfg, Turing.FinTM.bufferTape, Turing.MultiTapeTM.runFrom,
  Turing.MultiTapeTM.runFrom_succ_eq_step, Turing.MultiTapeTM.runFrom_zero,
  Turing.UniversalControl, Turing.universalEvalCfg,
  Turing.universalEval_step, Turing.universalFour,
  Turing.universalInterpreter}
\difficulty{4}
Suppose the table is $l \cdot 1^n \cdot 0 \cdot r$.  From the interpreter's
initial-state skipping phase (carrying a record index) with table cursor
$|l|$, exactly $n + 1$ steps move the table cursor past the unary field and
its terminating $0$ to position $|l| + n + 1$ and enter the record-group
phase with the carried index; the state tape, state cursor, and all other
data are unchanged.
\end{lemma}

\begin{lemma}[Exact-cost unary copy onto the state tape]
\lean{Turing.universal_unary_copy}
\label{Turing.universal_unary_copy}
\leanok
\uses{Turing.Cfg, Turing.FinTM.bufferTape, Turing.MultiTapeTM.runFrom,
  Turing.MultiTapeTM.runFrom_succ_eq_step, Turing.MultiTapeTM.runFrom_zero,
  Turing.UniversalControl, Turing.universalAdmin, Turing.universalEvalCfg,
  Turing.universalEval_step, Turing.universalFour,
  Turing.universalInterpreter, Turing.universalStateTape,
  Turing.universalStateTape_append}
\difficulty{4}
Let $q$ be an interpreter control state that, on every combination of reads
with table read $1$, stays in $q$, writes a $1$ at the state cursor, and
moves both the table and state cursors right; and on every combination with
table read $0$ passes to a state $q'$ with a fixed table movement
$d \in \{-1, 0, +1\}$ and no other change.  Suppose the table is
$l \cdot 1^n \cdot 0 \cdot r$, the state tape is the intact unary encoding of
an index $j$, the table cursor is $|l|$, and the state cursor is $j + 1$.
Then exactly $n + 1$ interpreter steps reach $q'$ with table cursor
$|l| + n + d$, state tape the intact unary encoding of $j + n$, and state
cursor $j + n + 1$; all other data are unchanged.  This single gadget serves
both initial-state extraction and live successor-state replacement.
\end{lemma}

\begin{lemma}[Uniqueness of the state-tape marker]
\lean{Turing.universalStateTape_marker}
\label{Turing.universalStateTape_marker}
\leanok
\uses{Turing.universalStateTape, Turing.universalStateWindow,
  Turing.universalStateWindow_zero}
\difficulty{3}
For every state index $n$, the intact unary state tape encoding $n$ reads the
marker bit $0$ at position $0$, and reads $0$ at no positive natural
position.
\end{lemma}

\begin{lemma}[Installing a permanent marker]
\lean{Turing.universal_install_marker}
\label{Turing.universal_install_marker}
\leanok
\uses{Turing.FinTM.bufferTape, Turing.FinTM.bufferTape_append}
\difficulty{1}
Updating the everywhere-blank two-way infinite tape by writing a bit $b$ at
position $0$ yields the buffer tape of the one-character string $b$: the tape
holding $b$ in cell $0$ and blanks elsewhere.
\end{lemma}

\begin{definition}[Interpreter entry configuration]
\lean{Turing.universalInterpreterInitial}
\label{Turing.universalInterpreterInitial}
\leanok
\uses{Turing.Cfg, Turing.FinTM.bufferTape, Turing.UniversalControl,
  Turing.universalFour}
The interpreter's entry configuration for a captured table $T$: control at
the interpreter's start state, physical input head at a given position $p$
(already parked at the start of the input suffix), first work tape holding
the buffer tape of $T$ with its head at $|T|$ (the right blank of the
table), the other three work tapes blank with heads at $0$, and no output
emitted.
\end{definition}

\begin{definition}[Initialization base configuration]
\lean{Turing.universalInterpreterBase}
\label{Turing.universalInterpreterBase}
\leanok
\uses{Turing.Cfg, Turing.FinTM.bufferTape, Turing.UniversalControl,
  Turing.universalFour}
The inactive data carried during interpreter initialization, packaged as a
base configuration: control at the interpreter's main dispatch state,
physical input head parked at a given position $p$, table and state tapes
blank with heads at $0$, simulated work tape blank with head at $0$, the
boundary-marker tape holding a single $1$ in cell $0$ with its head at
position $1$ (also correct for empty suffixes), and no output emitted.
\end{definition}

\begin{lemma}[The first interpreter step]
\lean{Turing.universalInterpreter_first}
\label{Turing.universalInterpreter_first}
\leanok
\uses{Turing.MultiTapeTM.step, Turing.moveInputPos_zero,
  Turing.universalEvalCfg, Turing.universalInterpreter,
  Turing.universalInterpreterBase, Turing.universalInterpreterInitial,
  Turing.universalStateTape, Turing.universal_install_marker}
\difficulty{3}
One step of the table interpreter from its entry configuration on a captured
table $T$ installs the two permanent markers --- the origin marker $0$ on the
state tape and the boundary bit $1$ on the marker tape --- and starts the
unconditional table rewind: the result is the evaluation configuration over
the initialization base whose control is the initial table-rewind phase, with
table cursor $|T| - 1$, the unary state tape encoding the index $0$, and
state cursor $1$.
\end{lemma}

\begin{definition}[Serialized transition record]
\lean{Turing.universalRecordBits}
\label{Turing.universalRecordBits}
\leanok
\uses{Turing.Action}
The bit-string serialization of one transition record of a one-work-tape
machine action over $\{0,1\}$: two bits encoding the input-head movement, two
bits encoding the optional work-tape write, two bits encoding the work-head
movement, two bits encoding the optional output emission, and finally the
optional successor state as a self-delimiting unary field --- a single $0$
when the action halts, and $1 \cdot 1^{q} \cdot 0$ when the successor state
has index $q$.
\end{definition}

\begin{definition}[Canonical record table of a coded machine]
\lean{Turing.universalRecords}
\label{Turing.universalRecords}
\leanok
\uses{Turing.CodeTM, Turing.universalRecordBits}
The canonical transition table of a coded machine $M$: the concatenation,
over every state of $M$ in order and over every pair of reads (input read,
work read) in the fixed serialization order blank, $0$, $1$, of the
serialized transition record of the action $M$ performs at that state and
read pair --- nine records per state, including all nine read pairs at every
live state.
\end{definition}

\begin{lemma}[Header shape of the canonical serialization]
\lean{Turing.universal_serialization_header}
\label{Turing.universal_serialization_header}
\leanok
\uses{Turing.CodeTM, Turing.CodeTM.serialize, Turing.pairEncode,
  Turing.universalRecords}
\difficulty{4}
For every coded machine $M$ with initial state index $q_0$ and state count
$m$, there is a bit string $\rho$ (the transition-record table) such that the
canonical serialization of $M$ equals the pair encoding of the binary digit
string of $m$ with the string $1^{q_0} \cdot 0 \cdot \rho$; that is, the
serialization begins with exactly its doubled count field and delimiter,
then the initial state in unary with its terminator, then the record table.
\end{lemma}

\begin{lemma}[Exact interpreter initialization]
\lean{Turing.universalInterpreter_initialize}
\label{Turing.universalInterpreter_initialize}
\leanok
\uses{Turing.MultiTapeTM.runFrom, Turing.MultiTapeTM.runFrom_add,
  Turing.MultiTapeTM.step, Turing.pairEncode, Turing.universalEvalCfg,
  Turing.universalInterpreter, Turing.universalInterpreterBase,
  Turing.universalInterpreterInitial, Turing.universalInterpreter_first,
  Turing.universalStateTape, Turing.universalStateTape_marker,
  Turing.universal_count_run, Turing.universal_pair_length,
  Turing.universal_state_rewind, Turing.universal_table_rewind,
  Turing.universal_unary_copy}
\difficulty{5}
Suppose the captured table $T$ is the pair encoding of a bit string $\beta$
(the count field) with a string of the form $1^n \cdot 0 \cdot \rho$ (the
initial state $n$ in unary, followed by an arbitrary remainder $\rho$).
Then from the interpreter's entry configuration with physical input head
parked at a position $p$, exactly $|T| + 2|\beta| + 2n + 7$ interpreter steps
reach the evaluation configuration over the initialization base at $p$ whose
control is the main dispatch state, with table cursor $2|\beta| + 2 + n + 1$
(just past the serialized header), the unary state tape encoding the index
$n$, and state cursor $1$.  No transition-table lookup is involved yet.
\end{lemma}

\begin{definition}[Interpreter checkpoint of a source configuration]
\lean{Turing.universalSimulationCfg}
\label{Turing.universalSimulationCfg}
\leanok
\uses{Turing.Cfg, Turing.CodeTM, Turing.CodeTM.serialize,
  Turing.FinTM.bufferTape, Turing.UniversalControl, Turing.pairEncode,
  Turing.universalFour, Turing.universalInputPos, Turing.universalStateTape}
For a coded machine $M$, a code string $\alpha$, a configuration of $M$'s
one-work-tape form on an input $x$, and a table-cursor value: the four-tape
interpreter configuration on the pair-encoded input
$\langle \alpha, x \rangle$ whose control is the main dispatch state when the
source is live and halted when the source has halted; whose physical input
head represents the source input head, offset past the encoded prefix; whose
table tape holds the buffer tape of the canonical serialization of $M$ at
the given cursor (the only administrative coordinate not determined by the
source); whose state tape holds the source's current state index in unary
with state cursor $1$; whose third work tape mirrors the source work tape
with the same head; whose marker tape holds a single $1$ at the origin with
its head at the source input position (so it reads $1$ exactly at the
simulated left boundary); and whose accumulated output equals the source
output.
\end{definition}

\begin{lemma}[Checkpoints preserve halting and output]
\lean{Turing.universalSimulation_fields}
\label{Turing.universalSimulation_fields}
\leanok
\uses{Turing.Cfg, Turing.CodeTM, Turing.universalSimulationCfg}
\difficulty{1}
The interpreter checkpoint representing a source configuration of a coded
machine is halted exactly when the source configuration is halted, and its
accumulated output equals the source's accumulated output.
\end{lemma}

\begin{lemma}[The capture endpoint is a right-block entry]
\lean{Turing.universalCaptured_right}
\label{Turing.universalCaptured_right}
\leanok
\uses{Turing.Cfg, Turing.FinTM, Turing.FinTM.bufferTape,
  Turing.FinTM.rightCfg, Turing.FinTM.tapeBlocks, Turing.MultiTapeTM,
  Turing.universalCaptureCfg, Turing.universalCapturedCfg,
  Turing.universalFour}
\difficulty{3}
Let $M$ be a finite-tape machine over $\{0,1\}$, $D$ a four-tape follower
machine, and consider a configuration of $M$ on an input $x$.  The endpoint
configuration of the capture stage --- which copies $M$'s accumulated output
onto a fresh tape and hands control to $D$ --- equals the right-block
embedding into the combined machine of the $D$-configuration with control at
$D$'s initial state, the same physical input head as the source, first work
tape holding the buffer tape of the captured output string with head at its
right blank, and three blank tapes with heads at $0$, where $M$'s work tapes
and head positions are retained unchanged as the inactive left block.
\end{lemma}

\begin{lemma}[Interpreter runs lift to the capture machine]
\lean{Turing.universalCapture_interpreter_run}
\label{Turing.universalCapture_interpreter_run}
\leanok
\uses{Turing.Cfg, Turing.FinTM, Turing.FinTM.rightCfg,
  Turing.FinTM.rightCfg_run, Turing.MultiTapeTM.runFrom,
  Turing.UniversalControl, Turing.universalCaptureTM,
  Turing.universalInterpreter, Turing.universalStartupBound}
\difficulty{1}
Let $M$ be a finite-tape machine over $\{0,1\}$.  For every four-tape
configuration of the table interpreter, every assignment of contents and head
positions to $M$'s work tapes, and every number of steps $t$: running the
combined capture machine (which pairs $M$ with the fixed table interpreter)
for $t$ steps from the right-block embedding of the interpreter configuration
equals the right-block embedding of the interpreter's own $t$-step run, with
the inactive tapes and heads unchanged.
\end{lemma}

\begin{definition}[The startup bound]
\lean{Turing.universalStartupBound}
\label{Turing.universalStartupBound}
\leanok
\uses{Turing.CodeTM.serialize, Turing.EffectiveMachineCode}
The startup time bound of the universal machine for a coding scheme on a code
string $\alpha$:
\[
  3|\alpha| + t_c(|\alpha|) + |s| + 2\ell + 2 q_0 + 12,
\]
where $t_c$ is the scheme's canonizer time bound, $s$ is the canonical
serialization of the machine decoded from $\alpha$, $\ell$ is the number of
binary digits of that machine's state count, and $q_0$ is its initial state
index.  It depends on $\alpha$ but not on the input suffix, and includes
complete interpreter setup.
\end{definition}

\begin{lemma}[Full prefix-start correspondence]
\lean{Turing.universal_initialized}
\label{Turing.universal_initialized}
\leanok
\uses{Turing.CodeTM.serialize, Turing.EffectiveMachineCode,
  Turing.FinTM.rightCfg, Turing.MultiTapeTM.initCfg,
  Turing.MultiTapeTM.runFrom, Turing.MultiTapeTM.runFrom_add,
  Turing.pairEncode, Turing.universalCanonTM, Turing.universalCanon_complete,
  Turing.universalCaptureTM, Turing.universalCapture_interpreter_run,
  Turing.universalCapture_start, Turing.universalCaptured_right,
  Turing.universalEvalCfg, Turing.universalFour, Turing.universalInterpreter,
  Turing.universalInterpreterBase, Turing.universalInterpreterInitial,
  Turing.universalInterpreter_initialize, Turing.universalSimulationCfg,
  Turing.universalStartupBound, Turing.universalTM,
  Turing.universal_serialization_header}
\difficulty{5}
Let $c$ be an effective machine-coding scheme, $\alpha$ a code string, and
$x$ an input string, and let $N$ denote the machine decoded from $\alpha$,
with $\ell$ the number of binary digits of its state count and $q_0$ its
initial state index.  Then there exist a time $t$ not exceeding the startup
bound for $\alpha$ (which is independent of $x$) and contents and head
positions for the inactive canonizer tapes such that running the complete
universal-machine candidate for $t$ steps from its initial configuration on
the pair encoding $\langle \alpha, x \rangle$ yields exactly the right-block
embedding, over those inactive tapes, of the interpreter checkpoint
representing the initial configuration of $N$ on $x$ with table cursor
$2\ell + 2 + q_0 + 1$ --- that is, with the canonical table captured, its
cursor parked just past the header, the initial state in unary, a blank
simulated work tape, and the virtual-input boundary marker installed.
\end{lemma}

\begin{lemma}[Cofinal physical run from a block simulation]
\lean{Turing.universal_block_run}
\label{Turing.universal_block_run}
\leanok
\uses{Turing.Cfg, Turing.MultiTapeTM, Turing.MultiTapeTM.initCfg,
  Turing.MultiTapeTM.runFrom, Turing.MultiTapeTM.runFrom_add,
  Turing.MultiTapeTM.runFrom_succ_eq_step', Turing.MultiTapeTM.step}
\difficulty{4}
Let $M$ and $U$ be multi-tape machines over $\{0,1\}$ on inputs $x$ and $y$
respectively, let $R$ be a relation between their configurations, and let
$S, B \in \mathbb{N}$.  Assume (startup) there is a $t \le S$ such that the
configuration of $U$ after $t$ steps from its initial configuration is
related to $M$'s initial configuration, and (step) whenever a source
configuration is related to a target configuration, there is a duration $d$
with $1 \le d \le B$ such that running $U$ for $d$ further steps yields a
configuration related to the source's successor.  Then for every $n$ there is
a time $t$ with $n \le t \le S + Bn$ such that the configuration of $U$ after
$t$ steps is related to the configuration of $M$ after $n$ steps.  The lower
bound $n \le t$ makes the run cofinal, without assuming that either machine
eventually halts.
\end{lemma}

\begin{lemma}[Evaluator clauses from a block simulation]
\lean{Turing.universal_from_blocks}
\label{Turing.universal_from_blocks}
\leanok
\uses{Turing.Cfg, Turing.CodeTM.toFinTM, Turing.EffectiveMachineCode,
  Turing.FinTM, Turing.FinTM.ComputesInTime,
  Turing.FinTM.ComputesInTime.mono, Turing.FinTM.computesInTime_iff,
  Turing.MultiTapeTM.initCfg, Turing.MultiTapeTM.runFrom,
  Turing.MultiTapeTM.step, Turing.pairEncode, Turing.universal_block_run}
\difficulty{5}
Let $c$ be an effective machine-coding scheme, $U$ a finite-tape machine over
$\{0,1\}$, and $S, B$ functions assigning a bound to each code string.  For
each code string $\alpha$ and input $x$, let a relation be given between
configurations of the machine decoded from $\alpha$ (in one-work-tape form)
on $x$ and configurations of $U$ on the pair encoding
$\langle \alpha, x \rangle$.  Assume: (startup) for all $\alpha, x$, within
some $t \le S(\alpha)$ steps $U$ reaches a configuration related to the
decoded machine's initial configuration; (step) every related pair advances
--- one source step against some $d$ target steps with
$1 \le d \le B(\alpha)$ --- to a related pair; (halting) related
configurations are halted simultaneously; and (output) related configurations
have equal accumulated output.  Then for every $\alpha$ there is a constant
$C$ such that for every input $x$: whenever the decoded machine computes an
output within $t$ steps on $x$, the machine $U$ computes the same output on
$\langle \alpha, x \rangle$ within $C(t+1)$ steps; and conversely, every
output $U$ computes on $\langle \alpha, x \rangle$ within some time is
computed by the decoded machine on $x$ within some time.
\end{lemma}

\begin{definition}[The eight fixed action bits]
\lean{Turing.universalActionBits}
\label{Turing.universalActionBits}
\leanok
\uses{Turing.Action}
The eight fixed bits of the serialized record of a one-work-tape machine
action, as a function on indices $0, \dots, 7$: bits $0, 1$ encode the
input-head movement, bits $2, 3$ the optional work-tape write, bits $4, 5$
the work-head movement, and bits $6, 7$ the optional output emission, each in
the two-bit encodings of the canonical record grammar.
\end{definition}

\begin{definition}[Unary size of the successor field]
\lean{Turing.universalNextOnes}
\label{Turing.universalNextOnes}
\leanok
The number of $1$s in the self-delimiting unary successor field of a
serialized record: $0$ when there is no successor state (the action halts),
and $q + 1$ when the successor state has index $q$ (one $1$ for the live flag
and $q$ for the index).
\end{definition}

\begin{lemma}[Shape of a serialized record]
\lean{Turing.universal_record_shape}
\label{Turing.universal_record_shape}
\leanok
\uses{Turing.Action, Turing.universalActionBits, Turing.universalNextOnes,
  Turing.universalRecordBits}
\difficulty{4}
The serialized record of any one-work-tape machine action decomposes as its
eight fixed action bits, followed by a block of $1$s whose length is the
unary size of the optional successor field (zero when the action halts, one
more than the successor index otherwise), followed by a single terminating
$0$.
\end{lemma}

\begin{lemma}[Exact-cost skip of the fixed fields]
\lean{Turing.universal_skip_fixed}
\label{Turing.universal_skip_fixed}
\leanok
\uses{Turing.Cfg, Turing.MultiTapeTM.runFrom,
  Turing.MultiTapeTM.runFrom_succ_eq_step, Turing.MultiTapeTM.runFrom_zero,
  Turing.UniversalControl, Turing.universalEvalCfg,
  Turing.universalEval_step, Turing.universalInterpreter}
\difficulty{4}
Consider the interpreter's fixed-field skipping phase, carrying a destination
and a count of remaining records, with field register $f \in \{0,\dots,7\}$
and table cursor $p$.  If $f + n = 7$, then exactly $n + 1$ interpreter steps
advance the table cursor to $p + n + 1$ and enter the unary-skipping phase
with the same destination and remaining count; no tape cell is inspected or
modified, and the state tape, state cursor, and all other data are unchanged.
\end{lemma}

\begin{definition}[Continuation after a skip request]
\lean{Turing.universalSkipDone}
\label{Turing.universalSkipDone}
\leanok
\uses{Turing.UniversalControl}
The interpreter control state entered after the last record of a skip
request: the record-group phase with index $i$ when the destination carries
an index $i$, and the action reader (starting at field $0$ with a cleared bit
buffer) when the destination is empty.
\end{definition}
